\section{Future Work}

This project work has the potential to be extended in several directions ranging from
immediate next steps to longer term research questions. The most obvious next step now would be 
to test the whole implementation on the real ARES accelerator rather than in simulation.  
A real machine could bring things that a simulator cannot completely capture, 
like measurement noise, small differences between cheetah model and actual machine etc. 
 Another interesting step 
for the future work could be to convert the code into a Python package. 
The system is currently tailored to ARES Experimental Area, with hardwired magnets names, lattice file, 
and misalignment values. But the idea of picking a cheetah lattice, tuning, learning and running BO 
itself is not specific to any one beamline, so packaging it up with a cleaner interface would 
make it much easier for others to try on their own accelerators.  \

Currently, each time the optimzer runs, it starts fresh, but in real accelerators the machine does not change drastically between the tuning
sessions. The learned misalignment values and GP hyperparameters could be used as warm start for the next,
potentially speeding up the convergence even further. So it would be worth having a look into it more. 
Some other research direction that could push this project further for future work could be 
tunable parameters increment. The ARES problem only has 5 tunable parameters that were the 5 magnets, 
but a real beamline can have dozens of them. At such scale, GP gets expensive and it gets 
harder for acquisition function to search, so combining Cheetah prior with dimensionality 
reduction techniques would be worth exploring. 


The current project uses objective function that combines beam position and size into a single MAE metric.
In some operational scenarios, it might be more natural to treat these as separate objectives 
(centering the beam and focusing the beam) and use multi-objective BO instead of combining everything to a single MAE,
so this could be another direction to experiment. Furthermore, on the modeling side, only the quadrupole misalignments were learned in this project, 
but there could also be other sources mismatched/unknown in the real machine, like incoming 
beam distribution may not be perfectly known, thus experiments can be conducted 
to learn other unknown parameters as well. Finally, the tests could also be conducted on more advanced beamlines with more 
complex objectives like the ones used in European XFEL. 

